% =====================================================================
% MockGen (standalone package) -- short conference paper, IEEEtran.
%
% WRITING RULES (enforced in review):
%  * Define every acronym at first use (OData, EDMX, CDS, MLP, ONNX,
%    VSIX, API, vCPU).
%  * Plain language. Short sentences. One idea per paragraph.
%  * Present tense, active voice, current-state facts only.
%  * Every measured number comes from numbers.tex.
%  * Every coverage number states the mode it was measured in.
%  * Label provenance (machine consensus, not human-verified) is stated
%    wherever label quality is discussed.
% =====================================================================
\documentclass[conference]{IEEEtran}
\usepackage[T1]{fontenc}
\usepackage[utf8]{inputenc}
\usepackage{booktabs}
\usepackage{graphicx}
\usepackage{microtype}
\usepackage{amsmath}
\usepackage{url}
\usepackage[hidelinks]{hyperref}
\Urlmuskip=0mu plus 1mu
\emergencystretch=3em

% =====================================================================
% numbers.tex -- single source of truth for every measured number cited
% by main.tex and conference.tex in this directory.
%
% RULES
%  * One \newcommand per metric; a source-path comment above each.
%  * Prose must use these macros -- a bare number in a .tex body is a
%    review defect.
%  * A value that cannot be traced to a file in this repository or to
%    the private evidence directory must not appear here.
%  * Values recorded 2026-09-21 (standalone package 0.1.0-dev.10).
% =====================================================================

% ---------------------------------------------------------------------
% Label production. Six adjudication rounds by a three-judge panel of
% distinct models over structural field metadata only.
% source: scripts/mockserver-data-generator-training/baselines/
%         2026-09-18-v3-development-evidence.json
% ---------------------------------------------------------------------
\newcommand{\LabelFields}{23{,}741}
\newcommand{\LabelRounds}{six}
\newcommand{\LabelJudges}{three}
\newcommand{\JudgeOne}{Claude Opus 4.8}
\newcommand{\JudgeTwo}{Claude Sonnet 4.5}
\newcommand{\JudgeThree}{GPT-5.6-terra}
\newcommand{\KappaLow}{0.72}
\newcommand{\KappaHigh}{0.87}
\newcommand{\KappaFinal}{0.81}

% ---------------------------------------------------------------------
% Partitions after the family-disjoint join.
% source: ~/mockgen-private/rows/join-summary.json
% ---------------------------------------------------------------------
\newcommand{\TrainRows}{12{,}375}
\newcommand{\CalibrationRows}{4{,}991}
\newcommand{\SealedRows}{4{,}317}
\newcommand{\TrainGroups}{119}
\newcommand{\CalibrationGroups}{39}
\newcommand{\SealedServices}{41}
\newcommand{\SealedDomains}{21}
\newcommand{\TrainedClasses}{63}
\newcommand{\RoutableRoles}{44}
\newcommand{\AuxiliaryRoles}{18}

% ---------------------------------------------------------------------
% Role head architecture and calibration.
% source: ~/mockgen-private/heads/v3-head.qualified.json
% ---------------------------------------------------------------------
\newcommand{\EmbeddingDim}{384}
\newcommand{\HiddenUnits}{256}
\newcommand{\TokenBudget}{64}
\newcommand{\CalibrationAccuracy}{0.835}
\newcommand{\ConformalCoverage}{0.900}
\newcommand{\GlobalThreshold}{0.54}
\newcommand{\RoleFloor}{5}
\newcommand{\FamilyFloor}{10}
\newcommand{\PrecisionTarget}{0.95}

% ---------------------------------------------------------------------
% Sealed evaluation of the shipped role head (service-disjoint).
% source: ~/mockgen-private/heads/sealed-report.dev10.json
% ---------------------------------------------------------------------
\newcommand{\SealedDecisions}{1{,}206}
\newcommand{\SealedCorrect}{1{,}160}
\newcommand{\SealedPrecision}{96.2}
\newcommand{\SupportedCorrect}{1{,}277}
\newcommand{\SupportedTotal}{2{,}026}
\newcommand{\SupportedRecall}{63.0}
\newcommand{\StatusCorrect}{27}
\newcommand{\StatusTotal}{44}
\newcommand{\StatusRawCorrect}{33}
\newcommand{\StatusServices}{22}
\newcommand{\StatusDomains}{10}
\newcommand{\CriticalFalsePositives}{27}
\newcommand{\CriticalRate}{2.2}
\newcommand{\LexicalAccepted}{512}
\newcommand{\MetadataAccepted}{389}

% ---------------------------------------------------------------------
% Serialization study: the same frozen encoder, two renderings of the
% same field context, linear head, identical partitions.
% source: session experiment logs, scratchpad serializer-experiment.mjs
% ---------------------------------------------------------------------
\newcommand{\DelimitedAccuracy}{0.63}
\newcommand{\SentenceAccuracy}{0.84}
\newcommand{\HeadNounPrecisionBefore}{45/56}
\newcommand{\HeadNounPrecisionAfter}{91/100}

% ---------------------------------------------------------------------
% Head-capacity study: linear versus one hidden layer, same embeddings.
% source: session experiment logs, scratchpad head_capacity.py
% ---------------------------------------------------------------------
\newcommand{\LinearCalAccuracy}{0.62}
\newcommand{\MlpCalAccuracy}{0.84}
\newcommand{\LinearRoutedShare}{29}
\newcommand{\MlpRoutedShare}{74}

% ---------------------------------------------------------------------
% Field-to-value relevance verifier.
% source: ~/mockgen-private/relevance3/heads/relevance-head.qualified.json
%         ~/mockgen-private/relevance3/heads/report.json
% ---------------------------------------------------------------------
\newcommand{\RelevancePairs}{3{,}346}
\newcommand{\RelevanceTrainGroups}{19}
\newcommand{\RelevanceSealedGroups}{19}
\newcommand{\RelevancePositives}{399}
\newcommand{\RelevancePositivesAccepted}{5}
\newcommand{\RelevanceNegatives}{704}
\newcommand{\RelevanceNegativesAccepted}{6}
\newcommand{\RelevanceNegativeRate}{0.9}
\newcommand{\RelevancePositiveRate}{1.3}
\newcommand{\RelevanceLooseNegativeRate}{14}

% ---------------------------------------------------------------------
% Packaged artifacts and release identities.
% source: scripts/mockserver-data-generator-training/baselines/
%         dev10-application-modeler-1.52.0.json
%         packages/mock-data-generator/resources/models/manifest.json
% ---------------------------------------------------------------------
\newcommand{\EncoderMiB}{21.9}
\newcommand{\RoleHeadMiB}{3.4}
\newcommand{\RelevanceHeadMiB}{0.1}
\newcommand{\VocabularyMiB}{0.2}
\newcommand{\GeneratorModelMiB}{157.3}
\newcommand{\PackageMB}{128}
\newcommand{\PackageLimitMB}{200}
\newcommand{\PackageVersion}{0.1.0-dev.10}
\newcommand{\ApiVersion}{2}
\newcommand{\VsixMB}{11.4}
\newcommand{\VsixVersion}{1.52.0}

% ---------------------------------------------------------------------
% Verification state of the shipped build.
% source: session test runs, 2026-09-20
% ---------------------------------------------------------------------
\newcommand{\PackageTests}{466}
\newcommand{\TrainingTests}{80}
\newcommand{\EditorTests}{50}
\newcommand{\LintErrors}{0}
\newcommand{\ReplayServices}{four}
\newcommand{\TravelAutoMs}{2{,}845}

% ---------------------------------------------------------------------
% Corpus benchmark over every checksum-verified service of the source
% registry whose format this package parses directly. Two runs of the same
% corpus: the recognition path (tiers 0, 1 and 3, the language model given a
% 1 ms budget so it never fires) and the whole pipeline with the language
% model at its packaged service budget. Field resolution is identical in
% both; only the timings differ.
% source: scripts/mockserver-data-generator-training/baselines/
%         2026-09-21-corpus-benchmark-recognition.json  (recognition path)
%         2026-09-21-corpus-benchmark-learned.json      (whole pipeline)
% ---------------------------------------------------------------------
\newcommand{\RegistryServices}{473}
\newcommand{\CorpusServices}{131}
\newcommand{\CorpusGenerated}{131}
\newcommand{\CorpusFailed}{0}
\newcommand{\CorpusEntities}{1{,}345}
\newcommand{\CorpusFields}{11{,}529}
\newcommand{\CorpusLargestService}{897}
\newcommand{\MetadataPct}{5.96}
\newcommand{\ClassifierPct}{29.01}
\newcommand{\LexicalPct}{11.85}
\newcommand{\AbstainedPct}{53.18}
\newcommand{\ProviderBoundPct}{44.79}
\newcommand{\CorpusMedianMs}{1{,}442}
\newcommand{\CorpusPNinetyMs}{2{,}929}
\newcommand{\CorpusTotalSec}{251}
\newcommand{\CorpusFullMedianMs}{3{,}808}
\newcommand{\CorpusFullPNinetyMs}{13{,}896}
\newcommand{\CorpusFullTotalSec}{820}

% ---------------------------------------------------------------------
% Same corpus with recognition disabled (deterministic mode): the
% contribution of the learned tier, and its effect on robustness.
% source: scripts/mockserver-data-generator-training/baselines/
%         2026-09-21-corpus-benchmark-deterministic.json
% ---------------------------------------------------------------------
\newcommand{\DetGenerated}{119}
\newcommand{\DetFailed}{12}
\newcommand{\DetFields}{9{,}928}
\newcommand{\DetMetadataPct}{5.96}
\newcommand{\DetLexicalPct}{17.94}
\newcommand{\DetAbstainedPct}{76.10}
\newcommand{\DetProviderBoundPct}{22.93}
\newcommand{\DetMedianMs}{107}

% ---------------------------------------------------------------------
% Scaling benchmark: deterministic pipeline, finance service, 17 entities.
% source: scripts/mockserver-data-generator-training/baselines/
%         2026-09-19-scaling-benchmark.json
% ---------------------------------------------------------------------
\newcommand{\ScaleEntities}{17}
\newcommand{\ScaleRowsMin}{17}
\newcommand{\ScaleRowsMax}{1{,}700}
\newcommand{\ScaleMsMin}{270}
\newcommand{\ScaleMsMax}{453}
\newcommand{\ScaleMsPerRowMax}{15.9}
\newcommand{\ScaleMsPerRowMin}{0.27}

% ---------------------------------------------------------------------
% Editor availability budgets.
% source: packages/application-modeler/ide-extension/src/mockgen/runtime.ts
%         (tools-suite-standalone-mockgen)
% ---------------------------------------------------------------------
\newcommand{\AvailabilityBudgetSec}{60}
\newcommand{\GenerationBudgetSec}{300}

\newcommand{\system}{MockGen}
\newcommand{\code}[1]{\texttt{#1}}

\begin{document}

\title{\system: Realistic Offline Mock Data for Enterprise Applications}

\author{\IEEEauthorblockN{Christos Koutsiaris}
\IEEEauthorblockA{\textit{SAP P\&E -- Cloud ERP -- UX Foundation}\\
Development Expert\\
christos.koutsiaris@sap.com}}

\maketitle

\begin{abstract}
Developers of enterprise applications need screens full of realistic
data long before any backend exists. Hand-written mock data is expensive
and goes stale; random generators produce values that look broken on
screen. \system{} generates complete, relationally consistent mock
datasets directly from an application's service definition, offline and
deterministically. Every column is resolved by a cascade: declared
catalogs, then a recognised semantic role, then a small local language
model, then a type-correct fallback that cannot decline. Recognition is
a learned classifier trained on \LabelFields{} fields labelled by a
panel of \LabelJudges{} large language models that see only structural
metadata. On \CorpusServices{} real service definitions the pipeline
generates \CorpusFields{} fields with \CorpusFailed{} failures and binds
\ProviderBoundPct\% of them to a semantic provider; with recognition
disabled the same corpus binds \DetProviderBoundPct\% and
\DetFailed{} services fail outright. On \SealedServices{} services held
out from training, the classifier's decisions are correct
\SealedPrecision\% of the time across \RoutableRoles{} roles. We also
report a verifier that did not work and what we shipped instead.
\end{abstract}

\begin{IEEEkeywords}
mock data, test data generation, service metadata, OData, offline
inference, semantic role classification, relational data generation
\end{IEEEkeywords}

\section{Introduction}
\label{sec:intro}

\IEEEPARstart{A}{frontend} developer starts from a service definition
--- a contract naming entities, columns, types, widths and relationships
--- and needs data on screen before any backend exists. Hand-written
mock data drifts out of sync with the schema. Random generators fill a
city column with random words, so every screen looks broken. Recording
real responses needs access, permissions and cleanup.

Enterprise schemas amplify this. A service declares dozens of entities
and hundreds of columns with exact widths; some columns are protocol
artifacts rather than business data; annotations encode meaning a
generator should honour. Errors are visible: a too-long value truncates,
a code that does not match its text reads as a bug, an empty child table
makes the application look defective.

\system{} generates a complete dataset from the service definition
alone: offline, deterministic, in seconds. Five rules govern it: any
application works cold; no screen is empty; data looks real at a glance;
no internal failure blanks the application; and the same input always
produces the same output. The architecture is a \emph{tier cascade}
(Fig.~\ref{fig:cascade}): the first tier that can produce a valid value
wins, and the last tier cannot decline.

\begin{figure}[t]
  \centering
  \includegraphics[width=0.74\columnwidth]{figures/out/fig-cascade.pdf}
  \caption{The cascade for one column. Declared catalogs beat
  recognition, recognition beats the language model, and the typed
  fallback cannot refuse. Takeaway: the never-empty promise is
  structural.}
  \label{fig:cascade}
\end{figure}

\textbf{Contributions.} (i) The four-tier cascade and the schema
representation behind it (Sections~\ref{sec:ir}--\ref{sec:tiers});
(ii) a learned role recogniser trained on machine-adjudicated labels,
with a routing policy that decides per role whether the model may speak
(Sections~\ref{sec:learning}--\ref{sec:routing}); (iii) a measured
characterisation of the shipped system, including a negative result
(Sections~\ref{sec:results}--\ref{sec:relevance}).

\section{Schema IR and Adapters}
\label{sec:ir}

Everything downstream reads one internal shape: a schema graph holding
entities, properties with type, width, precision, key and nullability
flags, declared labels, data elements and descriptions, relationships
with their column mappings, and annotation-derived links such as value
lists and text associations. Adapters translate each input format into
it, so the tiers never learn protocol dialects. Two ship today: the
metadata format (EDMX) of Open Data Protocol (OData) services, and Core
Data Services (CDS) documents.

\section{The Four Tiers}
\label{sec:tiers}

\textbf{T0, declared catalogs.} If the schema declares allowed values,
or the developer supplied authored rows, the value comes from there. A
value-list annotation points at a code entity; the generator produces
that code table first and draws from it, so a code column and its text
column agree by construction.

\textbf{T1, recognised roles.} Otherwise the system recognises what the
column means and draws from a curated bank. A role declares the types it
accepts, whether it may sit on a key, its provider, its validator, its
coherence group and its routing confidence. Roles group into families
--- status, identity, location, finance, temporal, narrative, contact,
organisation, measure --- and families decide both support floors and
error severity.

Three sources feed one arbitrator. Metadata rules fire when an
annotation or data element settles the question, and win outright. The
learned classifier handles the rest. Audited name-based rules remain
where a name is conclusive; each was checked against adjudicated labels,
and the eight that contradicted a label with a different positive role
were switched off. When classifier and lexical rule disagree, the
classifier wins only with a singleton prediction set and a stricter
confidence; otherwise both are discarded. Compatibility is enforced
last, so a role forbidden on keys never lands on one. Draft-administration
and action-control columns abstain unconditionally.

\textbf{T2, a budgeted local language model.} For visible text with no
recognised role, a small quantized model proposes values under a grammar
and a wall-clock budget, on two virtual central processing units
(vCPUs). When the budget expires the tier yields and the cascade
continues.

\textbf{T3, the guaranteed floor.} A typed generator emits values
respecting type, width, precision and nullability, unique where
uniqueness is required. It cannot decline, which is what turns ``no
empty screens'' into a property.

\textbf{Relational passes.} Keys are generated first and kept unique
within their entity. Child rows draw foreign keys from parents that were
actually generated, so references resolve. Coherence groups --- city
with country and postal code, amount with currency --- are generated
together, and a second pass recomputes derived display values. Datasets
are cached under a key including the schema, seed, row counts and every
model fingerprint, so a stale dataset is never served after a model
swap.

\section{Learning to Recognise a Column}
\label{sec:learning}

\subsection{Labels without human annotators}

We use a panel of \LabelJudges{} large language models from two vendors:
\JudgeOne{}, \JudgeTwo{} and \JudgeThree{}. Each sees one field's
structural metadata and a frozen guideline, and returns a registered
role or \code{unknown}. Judges are independent; a field is labelled when
at least two of three agree, and disagreements are excluded rather than
tie-broken. \LabelRounds{} rounds produced \LabelFields{} labels, with
Fleiss' kappa from \KappaHigh{} down to \KappaLow{} on a round aimed at
the hardest cases and back to \KappaFinal{} after the guideline was
clarified. Every label records machine panel consensus with
\code{humanVerified} false; we claim consistency and reproducibility,
not parity with expert labels.

Services group into families, and a family belongs to exactly one
partition: \TrainRows{} training rows from \TrainGroups{} groups,
\CalibrationRows{} calibration rows from \CalibrationGroups{}, and a
sealed set of \SealedRows{} fields across \SealedServices{} services.
Contexts with conflicting labels, and contexts appearing in two
partitions, are dropped rather than relabelled.

\subsection{Two findings that mattered}

The classifier embeds a text rendering of the field context with a
frozen sentence encoder (Fig.~\ref{fig:recognition}). The pilot rendered
a delimited record, \code{v3|BookingStatus\_code|string|...}, which is a
fine machine format and a poor input for an encoder trained on prose.
The shipped rendering is a sentence: \code{Booking Status code (string)
in Travel Booking; data element BOOK STATUS; a code field}. Identifiers
become words, structural facts are stated in words, and verbose material
that fills the \TokenBudget{}-word-piece budget without carrying role
evidence is dropped. With the same encoder and labels, calibration
accuracy rose from \DelimitedAccuracy{} to \SentenceAccuracy{}. The
final clause names the identifier's head noun, because mean pooling
cannot weigh the last word yet that word decides meaning; it raised
status routing precision from \HeadNounPrecisionBefore{} to
\HeadNounPrecisionAfter{}.

The second finding is capacity. One hidden layer of \HiddenUnits{}
rectified-linear units on the same embeddings raised calibration
accuracy from \LinearCalAccuracy{} to \MlpCalAccuracy{} and the routed
share of role-bearing calibration fields from \LinearRoutedShare\% to
\MlpRoutedShare\%. The shipped head is a multilayer perceptron (MLP) of
\EmbeddingDim{}$\times$\HiddenUnits{}$\times$\TrainedClasses{}, adding
\RoleHeadMiB{}\,MiB to the package.

\begin{figure}[t]
  \centering
  \includegraphics[width=0.70\columnwidth]{figures/out/fig-recognition.pdf}
  \caption{One column becoming a routed role: sentence rendering, frozen
  encoder, small head, three gates that can each fall back to
  abstention. Takeaway: routing is a decision about trust, not the
  top-scoring class.}
  \label{fig:recognition}
\end{figure}

\section{When the Classifier May Speak}
\label{sec:routing}

Three gates. \textbf{Support}: at least \RoleFloor{} correct calibration
examples for the role and \FamilyFloor{} for its family.
\textbf{Measured precision}: each role's threshold is raised --- never
lowered --- to the lowest confidence reaching \PrecisionTarget{}
precision on calibration with enough decisions to measure; a role that
cannot reach it does not route. \textbf{Conformal set size}: a decision
whose prediction set holds more than one class abstains, at a quantile
targeting \ConformalCoverage{} coverage.

Roles failing these gates stay as \emph{auxiliary} classes: trained,
never selected. When we instead removed weak classes from training,
status-text columns had nowhere to go and were misrouted as status
codes. Of \TrainedClasses{} classes, \RoutableRoles{} route and
\AuxiliaryRoles{} are auxiliary, so coverage grows with evidence rather
than with code changes.

\section{Results}
\label{sec:results}

\subsection{Corpus}

We run the whole pipeline over every checksum-verified service in the
source registry whose format the package parses directly:
\CorpusServices{} of \RegistryServices{} registered, covering
\CorpusEntities{} entities and \CorpusFields{} fields at two rows per
entity. All \CorpusGenerated{} generate, with \CorpusFailed{} failures.
\ProviderBoundPct\% of fields bind to a semantic provider: classifier
\ClassifierPct\%, metadata rules \MetadataPct\%, lexical rules
\LexicalPct\%.

Disabling recognition isolates the learned tier
(Fig.~\ref{fig:resolution}): binding falls to \DetProviderBoundPct\%
and, unexpectedly, \DetFailed{} services fail outright, because without
a recognised role some code columns cannot satisfy their own semantic
validation. Recognition is a robustness feature as well as a realism
one.

\begin{figure}[t]
  \centering
  \includegraphics[width=0.96\columnwidth]{figures/out/fig-resolution-bars.pdf}
  \caption{Field resolution across the \CorpusServices{}-service corpus
  with and without the learned tier. Takeaway: the classifier roughly
  doubles semantic binding, \DetProviderBoundPct\% to
  \ProviderBoundPct\%.}
  \label{fig:resolution}
\end{figure}

\subsection{Recognition quality on unseen services}

Corpus shares say how often the classifier speaks, not whether it is
right. The sealed set is family-disjoint from training and calibration
and is scored through the real runtime path. On \SealedServices{}
services and \SealedDomains{} domains never seen, the classifier makes
\SealedDecisions{} decisions of which \SealedCorrect{} are correct
(\SealedPrecision\%) across \RoutableRoles{} roles, with a critical
false-positive rate of \CriticalRate\%.

\begin{table}[t]
\centering
\caption{Sealed evaluation on \SealedServices{} unseen services
(\SealedRows{} fields). Precision counts the classifier's own decisions;
deterministic rules are listed separately. Takeaway: when the model
speaks it is right \SealedPrecision\% of the time.}
\label{tab:sealed}
\footnotesize
\begin{tabular}{@{}p{0.56\columnwidth}r@{}}
\toprule
Measure & Value \\
\midrule
Classifier routing decisions & \SealedDecisions{} \\
\quad correct & \SealedCorrect{} (\SealedPrecision\%) \\
\quad critical false positives & \CriticalFalsePositives{} (\CriticalRate\%) \\
Distinct roles routed & \RoutableRoles{} \\
Supported-field recall & \SupportedRecall\% \\
Status fields detected & \StatusCorrect{}/\StatusTotal{} \\
\quad classifier's top label correct & \StatusRawCorrect{}/\StatusTotal{} \\
Decisions from metadata rules & \MetadataAccepted{} \\
Decisions from lexical rules & \LexicalAccepted{} \\
\bottomrule
\end{tabular}
\end{table}

Two numbers belong together: status recall is
\StatusCorrect{}/\StatusTotal{} while the top label is correct on
\StatusRawCorrect{}/\StatusTotal{}. The first gap is the routing policy
declining what it cannot justify; the second is the model being wrong,
and only that closes with more labels. Residual errors concentrate on
one boundary --- a generic caption against a specific named thing ---
the same boundary where the judges themselves needed a guideline
revision.

\subsection{Speed and footprint}

Generation time is dominated by fixed costs, not dataset size
(Fig.~\ref{fig:scaling}): \ScaleRowsMin{} rows cost \ScaleMsMin{}\,ms
and \ScaleRowsMax{} rows \ScaleMsMax{}\,ms on a \ScaleEntities{}-entity
service, from \ScaleMsPerRowMax{} down to \ScaleMsPerRowMin{}\,ms per
row. Median service time is \DetMedianMs{}\,ms without recognition and
\CorpusMedianMs{}\,ms with it; both isolate recognition by giving the
language model a budget it cannot use, and letting it run at its packaged
budget raises the median to \CorpusFullMedianMs{}\,ms for the same
resolved fields. The reference application with every tier enabled
generates in \TravelAutoMs{}\,ms.

The published archive is \PackageMB{}\,MB against a
\PackageLimitMB{}\,MB registry limit: encoder \EncoderMiB{}\,MiB, role
head \RoleHeadMiB{}\,MiB, vocabulary \VocabularyMiB{}\,MiB, relevance
head \RelevanceHeadMiB{}\,MiB, generation model
\GeneratorModelMiB{}\,MiB. The editor extension (\VsixMB{}\,MB) carries
no models: it pins the package by version, application programming
interface (API) version and hashes, and resolves that pin once into a
shared verified cache. The build passes \PackageTests{} package tests,
\TrainingTests{} training-tool tests and \EditorTests{} editor tests
with \LintErrors{} lint errors.

\begin{figure}[t]
  \centering
  \includegraphics[width=0.96\columnwidth]{figures/out/fig-scaling.pdf}
  \caption{Deterministic-pipeline scaling on a \ScaleEntities{}-entity
  service. Takeaway: everything except the language model is effectively
  free at realistic dataset sizes.}
  \label{fig:scaling}
\end{figure}

\section{A Component That Did Not Work}
\label{sec:relevance}

Before a model-proposed display value is written, a verifier should
judge whether it suits the field; the intended gate was 80\% of valid
candidates accepted at no more than 1\% of wrong ones. We built
\RelevancePairs{} panel-judged pairs over \RelevanceTrainGroups{}
approved value sources with \RelevanceSealedGroups{} held out, and tried
a frozen-encoder logistic head, a fine-tuned bi-encoder and a
cross-encoder. None separated the classes: at the 1\% budget the
verifier accepts \RelevancePositiveRate\% of valid candidates, and
loosening it enough to be useful admits about
\RelevanceLooseNegativeRate\% of clearly-wrong values. Expanding the
corpus made the task harder, and the panel judged a sixth of our
supposed wrong values acceptable --- the task as posed was not well
defined.

We ship the strict point
(\RelevancePositivesAccepted{}/\RelevancePositives{} valid,
\RelevanceNegativesAccepted{}/\RelevanceNegatives{} wrong). Because that
rejects nearly everything and the pipeline treated ``unverified'' as
fatal, applications with synthetic display domains produced no data.
Generation now keeps the deterministic values, marks the resource
unverified and emits a diagnostic; the guardrail for a \emph{missing}
verifier is unchanged.

\section{Discussion and Limitations}
\label{sec:discussion}

The pilot recogniser routed nothing useful for three reasons, none of
them the model: the input was written for a machine rather than the
encoder, the head lacked capacity for the class count, and a registry
default threshold overrode the head's own calibration so decisions
between 0.6 and 0.9 never routed. Each was invisible in aggregate
accuracy and obvious once every sealed decision recorded its confidence,
threshold, prediction-set size and abstention reason.

Machine labelling bought scale --- \LabelFields{} labels for the cost of
compute --- but is also the ceiling: residual errors sit where
inter-judge agreement was lowest.

Limitations: labels are machine consensus, not human-verified; corpus
numbers count fields bound to a provider, not fields a human would judge
realistic; the corpus is \CorpusServices{} services of
\RegistryServices{} registered, and the sealed set, while
family-disjoint, comes from the same registry; supported recall is
\SupportedRecall\%; the verifier is nearly inactive; and gate thresholds
were revised once, with the owner's approval, when the sealed set grew
from 45 fields to \SealedRows{}, with both values recorded in the
release baseline.

\section{Conclusion}
\label{sec:conclusion}

\system{} turns a service definition into a complete, coherent,
deterministic mock dataset with no backend, no network and no human
labelling effort. On \CorpusServices{} real services it generates
\CorpusFields{} fields with \CorpusFailed{} failures and binds
\ProviderBoundPct\% to a semantic provider, against
\DetProviderBoundPct\% and \DetFailed{} failures without recognition; on
\SealedServices{} unseen services its learned decisions are correct
\SealedPrecision\% of the time across \RoutableRoles{} roles.

\section*{Glossary}

\begin{description}
  \item[API] Application programming interface.
  \item[CDS] Core Data Services, SAP's schema definition language.
  \item[Conformal prediction] A calibration method producing a set of
    plausible classes at a target coverage level.
  \item[EDMX] The XML document format for OData service metadata.
  \item[Fleiss' kappa] Agreement measure for more than two raters,
    corrected for chance.
  \item[MLP] Multilayer perceptron, a network with a hidden layer.
  \item[OData] Open Data Protocol, used by SAP Fiori applications.
  \item[ONNX] Open Neural Network Exchange, the packaged model format.
  \item[vCPU] Virtual central processing unit.
  \item[VSIX] The package format for a code editor extension.
\end{description}

\end{document}
